% rbd-support (base URLs set once; only ids and pull request numbers are given): \addjirainfo, \gls (monospace id,
% title in parentheses on first use), the plain acronym commands and the Jira glossary table.
\documentclass[aspectratio=169,10pt]{beamer}
\usetheme{joby}
\usepackage{rbd-support}
\setjirabaseurl{https://jira.example.com/browse/}
\setgithubpullbaseurl{https://git.example.com/pull/}
\addjirainfo{S4TC-1000}{FIRST-DESCRIPTION without a pull request}
\addjirainfo[42,43]{S4TC-1001}{SECOND-DESCRIPTION with a pull request}
\addjirainfo[7]{S4TC-1002}{THIRD-DESCRIPTION with one pull request}
\addjirainfo[5]{S4TC-1004}{FOURTH-DESCRIPTION used through gls}
\addjirainfo[9]{S4TC-1003}{UNREFERENCED-DESCRIPTION is never cited}
\newacronym{xx}{XX}{OTHER-ACRONYM-LONG}
\title{Support}
\author{Test Author}
\begin{document}
\begin{frame}{Citing tickets}
	Cites \acrshort{S4TC-1000} and \acrshort{S4TC-1001} in the text, and \acrshort{S4TC-1002}.
\end{frame}
\begin{frame}{Gls}
	First: \gls{S4TC-1004}. Again: \gls{S4TC-1004}.
\end{frame}
\begin{frame}{Short}
	Short: \acrshort{S4TC-1002}.
\end{frame}
\begin{frame}{Long}
	Long: \acrlong{S4TC-1002}.
\end{frame}
\begin{frame}{Other}
	Other: \gls{xx}, \acrfull{xx}.
\end{frame}
\begin{frame}{Full}
	Full: \acrfull{S4TC-1002}.
\end{frame}
\begin{frame}[allowframebreaks]{Tickets}
	\printjiraglossary
\end{frame}
\end{document}
